% Figure Legends (250 words per legend)
% - Should always have a short, explanatory title (as well as legend).
% - Legends must explain the figure fully without reference to the text.
% - If applicable, the original source of previously published material must be acknowledged in the Figure legend.
% - Figures must be cited in the main text (box figures must be cited in the box). 

\begin{figure}[p]
  \centering
\includegraphics[width=7cm]{figures/gaussian}
\caption{\textbf{Training and test inputs in the referential game of
    \cite{Bouchacourt:Baroni:2018}.} Two agents were trained to play
  the  game of \citeA{Lazaridou:etal:2017} (see
  Fig.~\ref{fig:referential-game}). During training, the agents were
  exposed to the same data as in the original study, that is, pairs of
  pictures of instances of about 500 distinct objects (top row). At
  test time, however, the agents were made to play the game with blobs
  of Gaussian noise (bottom row). They were able to
  communicate about them nearly as well as about the training
  pictures. This shows that the language emerging in this game does
  not involve ``words'' referring to generic concepts, but rather
  \emph{ad-hoc} signals, probably carrying comparative information
  about shallow visual properties of the images. Bottom row reproduced
  from \citeA{Bouchacourt:Baroni:2018} by permission.}
  \label{fig:gaussian}
\end{figure}
